to @@M0@@, whence the result, since the homomorphism is compatible with the filtrations.

\sgaStableTarget{sga6.e05.sec.6}\section*{6. Chern ring}

In this section, the base ring will be a binomial ring (2.7) @@M1@@, fixed once and for all. We shall use the category @@M2@@ (1.1), and therefore return to the notation of Section 1, reserving the notation @@M3@@ for the functor defined on @@M4@@ (1.1), and denoting by @@M5@@ its restriction to @@M6@@ (1.8).

\sgaStableTarget{sga6.e05.para.6-1}6.1. In 3.5 and 3.6 we defined morphisms of functors:
@@M7@@
the last being multiplication under the law @@M8@@ by @@M9@@, with @@M10@@.

By 1.9, we deduce morphisms:
@@M11@@
and it follows from the isobaricity properties of the polynomials @@M12@@, @@M13@@, @@M14@@ (3.5, 3.6, 2.9) that the image of the first and third functors is in @@M15@@, and that of the second in the subgroup of formal series in @@M16@@ whose coefficients in nonzero degree are in @@M17@@.

Composing the first and third with the canonical morphism: @@M18@@ (1.11), we obtain, by 1.12, two morphisms of functors on @@M19@@:
